\begin{solapartat}
\punts{0,2}
\[ f = \frac{c}{\lambda} = 10^{15}\ \mathrm{Hz} \]

\punts{0,3}
\[ E_{fot\acute{o}} = hf = 6,63\times 10^{-19}\ \mathrm{J} = 4,14\ \mathrm{eV}. \]

\punts{0,5}
\[ E_{C,\text{màx}} = hf - W_e = 2,78\times 10^{-19}\ \mathrm{J} = 1,74\ \mathrm{eV} \]

\punts{0,25}
\[ I = 2,00\times 10^{15}\,\frac{e}{s}\,\frac{1,602\times 10^{-19}\ \mathrm{C}}{1e} = 3,20\times 10^{-4}\ \mathrm{A} = 0,320\ \mathrm{mA} \]
(Com que no s'indica en quin sistema d'unitats s'ha de respondre, els resultats es poden donar tant en J com en eV.)
\end{solapartat}

\begin{solapartat}
\punts{0,25} La freqüència llindar correspon a la dels fotons que tenen una energia $W_e$:
\[ f = \frac{W_e}{h} = \frac{2,4\cdot 1,602\times 10^{-19}}{6,63\times 10^{-34}} = 5,80\times 10^{14}\ \mathrm{Hz} \]

\punts{0,1} La freqüència dels fotons del nou tub és: $f = \frac{c}{\lambda} = 8,57\times 10^{14}\ \mathrm{Hz}$

\punts{0,25} Per tant, també s'arrencaran electrons per efecte fotoelèctric amb el nou tub.

\punts{0,4} El nombre d'electrons arrencats depèn del nombre de fotons incidents, i com que el nombre de fotons és el mateix, el nombre d'electrons arrencats no varia.

\punts{0,25} L'única diferència serà que l'energia cinètica dels electrons emesos serà menor, atès que la freqüència (i, per tant, l'energia) dels fotons emesos pel nou tub és menor. Ara bé, el que és important és que els fotons generats pel nou tub tenen suficient energia per arrencar electrons, de manera que el dispositiu de seguretat segueix funcionant. En cas contrari, hauria estat necessari canviar aquest dispositiu.
\end{solapartat}
